\section{Incidence, order, perpendiculars, and affine geometry}
\label{sec:affine}
For $u=(u_1,u_2),v=(v_1,v_2)$ put
\[
 \dotp uv=u_1v_1+u_2v_2,\quad [u,v]=u_1v_2-u_2v_1,
 \quad N(u)=u_1^2+u_2^2,\quad \J u=(-u_2,u_1).
\]
A point is a pair of scalar codes. A valid line is $L(A,u)=\{A+tu:t\in\K\}$
with $N(u)>0$; braces abbreviate a formula, not a new set sort.
A segment is given by the same expression with $0\le t\le1$.
Parallel directions mean zero determinant; parallel supporting lines in
this directional sense may coincide. Strictly parallel lines are distinct
and disjoint. Whenever disjointness is intended it is stated explicitly.
The doubled signed area of $ABC$ is $[B-A,C-A]$.
The main article's analytic-geometry, parallel-postulate, Ceva, and
Menelaus links motivate this section \cite{wiki,euclid}.

\begin{theorem}[A line through two points and extension]\label{aff:line}
Two distinct points determine a unique line. Every nondegenerate segment can be extended
by a specified nonnegative multiple of its direction, and by any prescribed
nonnegative length along its ray.
\end{theorem}
\begin{proof}
For $A\ne B$ take $u=B-A$, so $N(u)>0$. The identity
\begin{equation}\label{eq:frame}
 N(u)v=\dotp uv\,u+[u,v]\J u
\end{equation}
is verified in its two coordinates by expansion. Thus $[u,v]=0$ exactly
when $v$ is a scalar multiple of $u$. A second line containing $A$ and
$B$ has a nonzero direction parallel to $u$, and hence consists of the
same points. The point $B+tu$, $t\ge0$, extends the segment. For a
specified additional length $l\ge0$, use $B+(l/\sqrt{N(u)})u$;
the denominator is positive and the new segment has squared length $l^2$.
\end{proof}

\begin{theorem}[Betweenness is a decidable polynomial predicate]
\label{aff:between}
Inclusive betweenness is precisely \eqref{eq:betweenpoly}. For $A\ne C$
its segment parameter is unique. If $A=C$, betweenness forces $X=A$.
\end{theorem}
\begin{proof}
Put $u=C-A$, $v=X-A$. If $u=0$, the dot condition is $N(v)\le0$,
so both coordinates of $v$ vanish. Otherwise $N(u)>0$ and
\eqref{eq:frame} turns the determinant condition into $v=tu$, with
$t=\dotp uv/N(u)$. The dot condition becomes $t(t-1)N(u)\le0$.
If $t<0$ or $t>1$, the two factors $t,t-1$ have the same nonzero sign,
contradicting it. Hence $0\le t\le1$. Conversely that interval condition
gives the determinant and dot inequalities directly. Uniqueness follows
by subtracting two parameters and using $N(u)>0$.
\end{proof}

\begin{theorem}[Complete line intersection test]\label{aff:intersection}
For $L(A,u),L(B,v)$ let $d=[u,v]$. If $d\ne0$, their unique intersection is
\[
 X=A+\frac{[B-A,v]}d\,u
   =B+\frac{[B-A,u]}d\,v.
\]
If $d=0$, the lines coincide when $[B-A,u]=0$, and are disjoint otherwise.
\end{theorem}
\begin{proof}
Write $w=B-A$. Expanding each coordinate gives
$[w,v]u-[w,u]v=[u,v]w$, proving both displayed expressions equal.
Conversely an intersection obeys $tu-sv=w$; taking determinants with
$v,u$ forces the two parameters. If $d=0$, formula~\eqref{eq:frame}
shows that the nonzero directions are scalar multiples. The displacement
$B-A$ is either a multiple too, giving the same line, or is not, excluding
every possible intersection. All three cases are decidable in $\K$.
\end{proof}

\begin{corollary}[Parallel through a point; Playfair and strong intersection]
\label{aff:parallel}
Through a point outside a line there is exactly one disjoint parallel line.
Two lines with nonparallel directions intersect at the explicitly computed
point of Theorem~\ref{aff:intersection}.
\end{corollary}
\begin{proof}
Through $B\notin L(A,u)$ use $L(B,u)$. The zero-determinant disjoint
case proves existence. Any other line through $B$ with determinant
nonzero against $u$ would meet $L(A,u)$; therefore its determinant is
zero and it is the same line through $B$. This conclusion uses the
decidable determinant of exact algebraic data. It is not a claim that
Playfair implies the same existence assertion in every intuitionistic
synthetic axiom system; that distinction is substantive \cite{beeson}.
\end{proof}

\begin{theorem}[Perpendicular foot and shortest distance]\label{aff:foot}
The unique foot of $P$ on $L(A,u)$ is
\[
 F=A+\frac{\dotp{P-A}u}{N(u)}u.
\]
For every $t\in\K$,
$N(P-(F+tu))=N(P-F)+t^2N(u)$.
Consequently the foot is the unique point of the line nearest to $P$.
The perpendicular through $P$ has direction $\J u$ and is unique.
\end{theorem}
\begin{proof}
Taking the dot product of $P-F$ with $u$ gives zero.
The stated square identity follows by expanding the norm and discarding
that cross term. Its last term is nonnegative, and is zero only at $t=0$.
Every vector perpendicular to $u$ is a multiple of $\J u$ by
\eqref{eq:frame}; this proves both the unique perpendicular and the unique
intersection which is its foot. The formula works when $P$ is already on
the line, so no separate incidence case is needed.
\end{proof}

\begin{theorem}[Perpendicular bisector as an exact locus]\label{aff:bisector}
For $A\ne B$, the points equidistant from $A,B$ are exactly the line
through $M=(A+B)/2$ perpendicular to $B-A$.
\end{theorem}
\begin{proof}
Set $u=B-A$. Direct expansion yields
$N(X-A)-N(X-B)=2\dotp{X-M}u$.
Equality of squared distances is therefore equivalent to perpendicularity,
and the positive square-root convention makes it equivalent to equality
of distances. The direction $\J u$ is nonzero.
\end{proof}

\begin{theorem}[Reflection in a line]\label{aff:reflection}
Let $F(P)$ be the perpendicular foot on a fixed line. The map
$R(P)=2F(P)-P$ is an involutive distance-preserving transformation, with
fixed points precisely that line.
\end{theorem}
\begin{proof}
Relative to $A$, decompose $P-A=\alpha u+\beta\J u$ by
\eqref{eq:frame}. The foot is $A+\alpha u$, and reflection replaces
$\beta$ by $-\beta$. Doing this twice is the identity. For two points
their difference decomposes in the same frame, and its squared norm is
$(\Delta\alpha)^2N(u)+(\Delta\beta)^2N(u)$, unchanged by the sign reversal.
A fixed point has $2\beta\J u=0$, hence $\beta=0$ and lies on the line.
\end{proof}

\begin{theorem}[Barycentric coordinates and oriented subareas]
\label{aff:barycentric}
For a noncollinear triangle $ABC$, every $X$ has unique coordinates
$X=\alpha A+\beta B+\gamma C$ with $\alpha+\beta+\gamma=1$.
With $D=[B-A,C-A]\ne0$,
\[
 \beta=\frac{[X-A,C-A]}D,\qquad
 \gamma=\frac{[B-A,X-A]}D,\qquad \alpha=1-\beta-\gamma.
\]
The three oriented subareas satisfy
$[B-X,C-X]=\alpha D$, $[C-X,A-X]=\beta D$,
$[A-X,B-X]=\gamma D$. Nonnegative coordinates characterize the closed
filled triangle; positive coordinates characterize its interior.
\end{theorem}
\begin{proof}
Subtract $A$ and solve the two-by-two linear system in the basis
$B-A,C-A$, using the determinant calculation in
Theorem~\ref{aff:intersection}. This gives existence and uniqueness.
Substitute $X-A=\beta(B-A)+\gamma(C-A)$ into the three determinants;
bilinearity gives the formulas. A convex combination has nonnegative
coefficients summing to one by definition, and uniqueness gives the converse.
The strict case lies on the inside side of each oriented edge because
all three determinant signs equal the sign of $D$; a zero coefficient is
on the corresponding side. These are finite sign tests.
\end{proof}

\begin{theorem}[Inner Pasch as a positive explicit intersection]
\label{aff:pasch}
Let $P=(1-r)A+rC$ and $Q=(1-s)B+sC$, with $0<r,s<1$ and $ABC$
noncollinear. The open segments $BP$ and $AQ$ intersect.
\end{theorem}
\begin{proof}
Put $H=r+s-rs>0$. Their common point has barycentric coordinates
\[
 X=\frac{s(1-r)A+r(1-s)B+rsC}{H}.
\]
The numerators are positive and sum to $H$. In $BP$ its parameter
from $P$ to $B$ is $r(1-s)/H$, strictly between zero and one; the
remaining coefficients equal $(1-r)$ and $r$ times its complement.
In $AQ$ its parameter from $A$ to $Q$ is $r/H$, also in $(0,1)$
since $H-r=s(1-r)>0$. Expanding the two segment formulas gives exactly
the displayed point. There is no order case split or unbounded search.
\end{proof}

\begin{theorem}[Intercept theorem and midpoint theorem]\label{aff:intercept}
Let $ABC$ be noncollinear, $X=A+t(B-A)$, $Y=A+s(C-A)$, and suppose
$X\ne Y$. Then $XY$ is parallel to $BC$ exactly when $t=s$.
In that case $Y-X=t(C-B)$, so lengths on the parallel segments have
ratio $|t|$. In particular the midpoint segment is parallel to the
third side and has half its length.
\end{theorem}
\begin{proof}
Translate $A$ to zero and write $B=u,C=v$. Bilinearity gives
$[sv-tu,v-u]=(s-t)[u,v]$. Since $[u,v]\ne0$, the determinant
vanishes exactly when $s=t$. Substitution gives the vector identity.
Taking norms proves the ratio and putting $t=1/2$ proves the midpoint case.
\end{proof}

\begin{theorem}[Ceva's theorem for internal cevians]\label{aff:ceva}
Let $D=(1-d)B+dC$, $E=(1-e)C+eA$, $F=(1-f)A+fB$, where
$0<d,e,f<1$. The lines $AD,BE,CF$ concur if and only if
\[
 \frac d{1-d}\frac e{1-e}\frac f{1-f}=1.
\]
These are the directed side ratios $BD/DC$, $CE/EA$, $AF/FB$.
\end{theorem}
\begin{proof}
Write the concurrency point in barycentric coordinates $(x,y,z)$.
The three line equations are
\[
 dy=(1-d)z,\qquad (1-e)x=ez,\qquad fx=(1-f)y.
\]
The first two imply that its coordinates are proportional to
$(de,(1-d)(1-e),d(1-e))$, whose sum is positive. Substitution in the
third gives $def=(1-d)(1-e)(1-f)$, equivalent to the displayed ratio.
Conversely, under that equality normalize these three positive coordinates
by their sum. They solve all three equations, constructing the common
interior point. The positivity excludes a point at infinity; directed
external variants require retaining that qualification.
\end{proof}

\begin{theorem}[Menelaus' theorem with directed ratios]\label{aff:menelaus}
With $D,E,F$ as in Theorem~\ref{aff:ceva}, but now arbitrary scalar
$d,e,f$ different from $0,1$, the points $D,E,F$ are collinear if and
only if
\[
 \frac d{1-d}\frac e{1-e}\frac f{1-f}=-1.
\]
\end{theorem}
\begin{proof}
In the affine barycentric chart the determinant of their coordinate rows is
\[
 \det\begin{pmatrix}0&1-d&d\\ e&0&1-e\\1-f&f&0\end{pmatrix}
  =(1-d)(1-e)(1-f)+def.
\]
The coordinate change has nonzero determinant, so collinearity is
equivalent to this expression being zero. Divide by the three nonzero
denominators to obtain the assertion. The signs automatically distinguish
external division from internal division; three internal division points
cannot satisfy this condition.
\end{proof}

\begin{theorem}[Affine maps and area scaling]\label{aff:maps}
For an invertible $2\times2$ matrix $T$, the map $X\mapsto a+TX$
preserves collinearity, segment parameters, parallelism, and concurrency.
It multiplies signed doubled areas by $\det T$, and unsigned areas by
$|\det T|$. Similarities of ratio $\lambda>0$ therefore multiply
areas by $\lambda^2$.
\end{theorem}
\begin{proof}
An affine combination with coefficient sum one commutes with the map.
This proves the incidence and parameter claims, and invertibility gives
their converses. Expanding the two-by-two determinant gives
$[Tu,Tv]=\det(T)[u,v]$. For a similarity $T=\lambda Q$ with
$Q^\mathsf TQ=I$, taking determinants gives $(\det Q)^2=1$,
so $|\det T|=\lambda^2$. No statement about arbitrary measurable sets
is hidden here: polygonal areas are finite sums, dealt with below.
\end{proof}

\begin{theorem}[A finite homogeneous Desargues theorem]\label{aff:desargues}
In homogeneous coordinates in $\K^3$, suppose two nondegenerate triangles
have representatives $A_i$ and $B_i=A_i+t_iO$ for $i=1,2,3$, with
each $t_i\ne0$. Assume each pair of corresponding side lines is distinct.
Their three intersection points are collinear, allowing a point at infinity.
For intersections with nonzero last coordinate this gives the usual
Euclidean-plane conclusion after division by that coordinate.
\end{theorem}
\begin{proof}
A projective point is a nonzero vector modulo nonzero scalar multiple;
a line is a two-dimensional span. Put
$P_{ij}=t_jA_i-t_iA_j$. It is nonzero because $A_i,A_j$ are independent.
The equality
$t_jB_i-t_iB_j=t_jA_i-t_iA_j$ places it in both side spans. Their
distinctness makes this their intersection. Finally
\[
 t_3P_{12}-t_2P_{13}+t_1P_{23}=0
\]
by cancellation of the three $A_i$ coefficients. The coefficient of
$P_{23}$ is nonzero. Moreover $P_{12},P_{13}$ are independent: in the
$A_1,A_2,A_3$ basis, a relation $aP_{12}+bP_{13}=0$ has second and
third coordinates $-at_1,-bt_1$, forcing $a=b=0$. Thus all three
points belong to the same genuine two-dimensional span.
All claims are finite linear algebra. The explicit representative
hypothesis states the nonexceptional perspective configuration being used;
it does not silently assert that every degenerate projective picture has
three defined intersections.
\end{proof}

\paragraph{iPRA audit.}
Every denominator here is $N(u)$, a nonzero determinant, a positive
barycentric sum, or a stated nonzero directed ratio denominator. Expand
and multiply through these denominators to obtain integer polynomial
identities and sign implications. Thus the proofs are fixed finite
ordered-field arguments to which Section~\ref{sec:foundation} applies.
No induction on geometric objects or quantification over arbitrary sets
occurs.
